\section{系统架构与协同接口}
\subsection{分层控制}
人形机器人负责物体搬运与领行；TurtleBot负责各自的目标观测、定位和运动控制。不同机器人的执行器由各自进程独占，感知线程只提供观测，避免多个程序同时写舵机或底盘速度。图\ref{fig:architecture}描述方式二的完整数据链；任务心跳是可选适配接口，不代表所有人形入口已经发送心跳。
\begin{figure}[H]\centering
\begin{tikzpicture}[flow/.style={-{Stealth[length=2mm]},thick,draw=Navy!85}, stage/.style={draw=Navy!45,rounded corners=2pt,fill=Navy!3,align=center,text width=35mm,minimum height=17mm,font=\small,inner sep=3mm}]
\node[stage] (v) at (0,0) {TonyPi 相机\\固定标记／道路观测};
\node[stage] (s) at (5,0) {路线状态机\\抓取、直行、转弯、放置};
\node[stage] (a) at (10,0) {动作组执行\\抱物上肢固定};
\draw[flow] (v)--(s);\draw[flow] (s)--(a);
\node[stage] (m) at (0,-2.7) {TurtleBot 相机\\背部 AR 目标};
\node[stage] (t) at (5,-2.7) {按时间戳变换至 map\\历史轨迹与回退点};
\node[stage] (n) at (10,-2.7) {\texttt{move\_base} / DWA\\地图与实时扫描};
\draw[flow] (m)--(t);\draw[flow] (t)--(n);
\node[stage] (g) at (10,-5.1) {许可与速度门控\\输出至底盘};
\draw[flow] (n)--(g);
\node[draw=Teal!55,rounded corners=2pt,fill=Teal!4,align=left,text width=85mm,inner sep=3mm,font=\footnotesize] (h) at (2.5,-5.1) {\textbf{\color{Navy}任务状态接口}\\[1.5mm]通用 \texttt{competition}：已支持心跳桥。\\最新 \texttt{marker\_route}：需适配任务状态输出。};
\draw[flow,dashed,draw=Teal] (h.east)--(g.west);
\end{tikzpicture}

\caption{人形机器人与方式二跟随的数据链}\label{fig:architecture}
\end{figure}

图中两条实线链分别对应人形动作控制和移动底盘导航，二者通过任务语义协同。底部虚线表示尚需适配的任务许可接口：它必须把领航者的运行状态显式传给跟随门控，不能由相机“仍能看见目标”代替。因而感知有效、任务允许运动和底盘速度输出应分别核对。


\subsection{软件平台与版本边界}
本稿主体最初依据\code{7bbf588}分析，新增视频后对GitHub主分支\code{43a7771}进行了针对性核对。TonyPi路线入口、前车位姿转发、地图跟随器和轨迹核心逻辑保持一致；当前\code{follow_map.launch}将速度门控指向\code{ar_marker_following}功能包，而非旧快照中的\code{multi_robot_following}。下述架构功能不变，但复现命令应使用当前包归属。复核文件保存在\code{code_snapshot_current}，原始分析快照保留在\code{code_snapshot}，以免将不同版本证据混写\cite{repo-current}。

TonyPi 使用 Python 3、OpenCV、NumPy、PyYAML及设备原装舵机SDK；不同树莓派镜像的串口、SDK和相机接口存在差异，应使用实际部署配置。TurtleBot课程虚拟机为 Ubuntu 16.04、ROS Kinetic、Python 2.7，功能包放置于 \code{~/catkin_ws/src}，不能直接用 Python 3 替换 ROS 节点解释器。

\begin{table}[H]\centering\small
\caption{关键接口及约束}
\begin{tabularx}{\linewidth}{p{37mm}p{40mm}Y}\toprule
接口 & 内容 & 使用约束\\\midrule
\code{/ar_pose_marker} & 带ID与时间戳的标记位姿 & ALVAR标记；尺寸、相机内参和坐标系须一致\\
\code{target_pose} & 目标在地图中的位姿 & 按观测时刻变换，保留原时间戳\\
\code{/scan} & 障碍扫描 & 实际可能来自Kinect深度转扫描；检查高度覆盖与时效\\
\code{map -> odom} & 全局定位变换 & 建图用GMapping；固定地图定位用AMCL，不能同时发布\\
\code{/competition/mission} & 搬运任务状态 & 需持续更新；开环和最新标记入口不能默认视为已接入\\
\code{/cmd_vel_mux/input/navi} & 方式二底盘导航速度 & 由独立门控输出；不与键盘遥控争用\\\bottomrule
\end{tabularx}\end{table}

这些接口共同要求时间、坐标和发布权限一致。消息能够接通只是第一步，还需核对观测采样时刻、TF 链以及最终速度发布者，才能判断系统是否按图示链路运行。


\subsection{拓扑与跟随方式的区分}
“TonyPi到首车再到后车”的队列拓扑，与“标记/地图/RGB-D/前车轨迹”的跟随算法是两个维度。当前方式二默认直接观察 TonyPi 背部 AR 标记，再在地图中导航；也支持输入前车地图位姿。实验时应固定一个目标来源并记录，不能把两个适配器同时发布到同一个目标话题。

多机应使用唯一节点名、命名空间与TF名称，并确保时间与地图原点一致。虚拟机 NAT 地址可供宿主机 SSH 访问，不保证机器人能反向访问 ROS；联调需验证共同局域网、\code{ROS_MASTER_URI}与\code{ROS_IP}，报告不收录密码等登录凭据。
